% Bìa báo cáo nhóm dựa trên khung đôi và logo trong template thesis UET.
\begin{titlepage}
\thispagestyle{empty}
\begin{tikzpicture}[overlay,remember picture]
 \draw[line width=3pt]
  ($(current page.north west)+(25mm,-25mm)$)
  rectangle ($(current page.south east)+(-15mm,25mm)$);
 \draw[line width=1pt]
  ($(current page.north west)+(26.5mm,-26.5mm)$)
  rectangle ($(current page.south east)+(-16.5mm,26.5mm)$);
\end{tikzpicture}
\begin{center}
\vspace*{3mm}
{\fontsize{13}{17}\selectfont\bfseries
 ĐẠI HỌC QUỐC GIA HÀ NỘI\\
 TRƯỜNG ĐẠI HỌC CÔNG NGHỆ\par}
\vspace{10mm}
\includegraphics[width=30mm]{figures/uet.jpg}\par
\vspace{9mm}
{\fontsize{16}{20}\selectfont\bfseries BÁO CÁO ĐỀ TÀI\par}
\vspace{13mm}
{\fontsize{21}{28}\selectfont\bfseries
 PE-GH: TRUY XUẤT LAI\par
 CÓ CHỌN LỌC TRƯỚC BƯỚC\par
 MÃ HÓA TRUY VẤN\par}
\vspace{7mm}
{\fontsize{13}{17}\selectfont\ReportSubtitle\par}
\vspace{17mm}
{\bfseries Nhóm thực hiện\par}
\vspace{4mm}
\begin{tabular}{ll}
 \textbf{Họ và tên} & \textbf{Mã số sinh viên}\\[3mm]
 \ReportStudentRows
\end{tabular}
\vfill
{\bfseries\MakeUppercase{\ReportPlace}, \ReportYear\par}
\vspace{8mm}
\end{center}
\end{titlepage}

